\section{Literature survey}\label{sec:literature}

Five recent works shape this project. Two study homecoming travel in China, one of them against the census of where migrants come from. One describes how holidays reorganise a national mobility network. Two study seasonal movement and return tourism in rural Spain. For each, I say what it did, what I take from it and where it falls short for the present question.

\subsection{Holiday travel against the census: the Chinese Spring Festival}

\citet{li2016spring} built networks of travel between China's 31 provinces during the 2015 Spring Festival, the largest annual return home in the world, from three online sources: location data from Baidu and Tencent, and train-ticket searches on Qihoo. They compared them with the 2010 census network of migrants, which links each migrant's registered province of origin to the province where they live. The two networks were closely aligned: the correlation between census and Baidu flows was 0.88, and 0.81 on a logarithmic scale, tested with the quadratic assignment procedure (QAP), which permutes the nodes of one network to judge whether two networks on the same nodes are more similar than chance.

This is the closest precedent to my question, and I borrow its central idea: a network of where people come from, set against a network of holiday travel on the same nodes. Its limits are also what this project tries to address. First, with 31 provinces the network is almost complete, so its structure (degrees, mixing, rich club) cannot be studied; only the weights can. Second, both networks fall with distance, and the correlation is not compared with what distance alone would give, so part of the alignment may be geography. Third, the census describes recent migrant workers, not a migration that happened two generations ago. Fourth, Baidu and Tencent publish only the ten largest flows of each city, which biases the network towards big links. The authors use the census mainly to check the quality of the new data, not as a hypothesis about why people travel.

\subsection{How distance changes on homecoming holidays}

\citet{luan2024different} modelled flows between Chinese cities during the same festival with Tencent data and gradient-boosted trees. They found that the usual fall of flows with distance does not hold on homecoming trips. Depending on the province, the relation shows a plateau with almost no decay, a drop, or a drop followed by a rebound at long distances.

The lesson for this project is about the null model. If homecoming changes how flows depend on distance, a distance model fitted on ordinary days cannot serve as the baseline for holiday nights. I therefore never compare a holiday night with a model of ordinary travel: every model is fitted on each night separately. In the roots test, two cities of the same province are compared against the same destination province on the same night, so what the two provinces share cancels out, and the distance from each city is modelled for that night (Section~\ref{sec:method}). The study's own weakness is that it never measures where travellers come from: the destination is called home because the trip happens at a homecoming holiday. Gradient-boosted trees describe the shape of the relation, but do not say why it differs between provinces.

\subsection{Holidays reorganise a national network}

\citet{wesolowski2017seasonal} used mobile-phone records from Kenya, Namibia and Pakistan to measure travel over the year. National holidays change the reach of travel, routes gain and lose importance from one season to another, and travel from rural to urban areas rises after the holidays. They show that these seasonal swings change how fast a disease spreads through the country.

From this work I take two choices. A holiday is treated as a different network on the same nodes, not as more of the same travel; and the direction of a link matters, since sending and receiving zones play different roles. Its limits for my question are that travel is not linked to where people were born, each country is covered by one phone operator, and the analysis works with flows between regions without measuring network structure. The setting is also very different from Spain, so what transfers is the method rather than the result.

\subsection{Seasonal population in Spain from phone data}

\citet{garzon2026heritage} used the aggregated phone indicators of the Spanish statistical office for 3,214 mobility areas on four dates of 2021. For each area they computed how many residents stay, how many leave and how many visitors arrive, and grouped the areas into three profiles: areas whose residents disperse (873), broadly connected areas (992) and areas that receive people in summer (1,344). The number of protected heritage sites in an area explained an additional 3.6--5.0\,\% of the variation between areas.

This is the most recent study of seasonal redistribution with Spanish phone data, and its summer-receiving interior areas correspond to the emptied villages of Section~\ref{sec:results}. I take from it the caution that phone data do not record the purpose of a trip, and treat single dates as snapshots that must be checked against similar dates. Its limits are that everything is measured per area: where visitors come from is reduced to a single diversity score, so nothing can be said about links between particular places. It has no information on birthplace, which is the variable my question needs. It also uses 2021, a pandemic year, and only four dates.

\subsection{Return tourism in rural Spain}

\citet{garciagonzalez2024retorno} measured \emph{return tourism} in Castilla-La Mancha, one of the regions that lost most people in the exodus: people who spend their free time in the municipality they or their families come from. They used the experimental phone statistics of the Spanish statistical office to estimate how many visitors each municipality receives, compared these numbers with the municipality's registered population and with its regulated tourist accommodation, and mapped the result. They read the visitors that regulated accommodation cannot explain as return tourism, a part of tourism that the official statistics do not see.

This is the closest Spanish precedent to my question, and it supports its premise: the return to the village is large enough to show in phone data. I take from it the idea of comparing what a place receives with what its population and tourist supply would predict. Its limits are the ones this project addresses. Return is inferred from where people go, not measured: the study does not know where the visitors were born or where they live, so a visitor to a village and an emigrant coming home look the same. It covers one region, describes places one by one rather than the links between them, and has no null model for distance or size.

\subsection{What is left open}

Together, these works show that holidays reorganise mobility networks \citep{wesolowski2017seasonal}, that homecoming travel follows the map of migration \citep{li2016spring} and does not decay with distance in the usual way \citep{luan2024different}, that rural Spain fills up in summer \citep{garzon2026heritage}, and that much of this is return tourism \citep{garciagonzalez2024retorno}. None of them links the holiday network of a country to the birthplaces of its residents at a scale where network structure can be measured, with controls for distance and for what the degrees alone would produce. In Spain, return to the village has been measured place by place and region by region, but not as a network that links the cities where people live to the provinces where they were born. This report does that, for 2,435 zones and with the comparisons in Section~\ref{sec:method}.
